\documentclass{article}

\usepackage{amsmath,amssymb}
\usepackage{xurl}
\usepackage[hidelinks]{hyperref}
\setlength{\emergencystretch}{2em}

% Shared commands for notes in docs/paper-gaps/.

\DeclareUnicodeCharacter{2080}{\ifmmode{}_{0}\else\textsubscript{0}\fi}
\DeclareUnicodeCharacter{2081}{\ifmmode{}_{1}\else\textsubscript{1}\fi}
\DeclareUnicodeCharacter{2082}{\ifmmode{}_{2}\else\textsubscript{2}\fi}
\DeclareUnicodeCharacter{2099}{\ifmmode{}_{n}\else\textsubscript{n}\fi}

\newcommand{\Lean}{\textsc{Lean}}
\ExplSyntaxOn
\NewDocumentCommand{\leanid}{m}{
  \ifmmode
    \textsf{\detokenize{#1}}
  \else
    \group_begin:
    \urlstyle{sf}
    \tl_set:Nn \l_tmpa_tl {#1}
    \tl_replace_all:Nnn \l_tmpa_tl {\_} {_}
    \exp_args:NV \path \l_tmpa_tl
    \group_end:
  \fi
}
\ExplSyntaxOff
\providecommand{\tr}{\operatorname{tr}}
\newcommand{\tnleanIssueBase}{https://github.com/LionSR/TNLean/issues/}
\newcommand{\tnleanPrBase}{https://github.com/LionSR/TNLean/pull/}
\newcommand{\tnleanDocBase}{https://sirui-lu.com/TNLean/docs/}
\newcommand{\ghissue}[1]{\href{\tnleanIssueBase#1}{GitHub issue \##1}}
\newcommand{\ghpr}[1]{\href{\tnleanPrBase#1}{PR \##1}}
\newcommand{\doclink}[1]{\href{\tnleanDocBase#1.html}{\path{#1}}}

% Dirac notation and the identity operator, as in blueprint/src/macros/common.tex.
% \providecommand keeps note-local definitions authoritative.
\usepackage{dsfont}
\providecommand{\ket}[1]{|#1\rangle}
\providecommand{\bra}[1]{\langle #1|}
\providecommand{\braket}[2]{\langle #1 | #2 \rangle}
\providecommand{\ketbra}[2]{|#1\rangle\!\langle #2|}
\providecommand{\Id}{\mathds{1}}
\providecommand{\one}{\mathds{1}}
\providecommand{\C}{\mathbb{C}}
\providecommand{\MN}[1]{M_{#1}(\C)}
\providecommand{\MD}{M_D(\C)}

% External referencing for paper sources.  The build helper writes sanitized
% auxiliary files under build/paper-aux/ and excludes duplicated source labels.
\usepackage{xr}

% Import a generated paper auxiliary file with a caller-supplied label prefix.
% The candidate paths support builds from docs/paper-gaps/, the repository root,
% and build/paper-aux/ itself.
\newcommand{\importpaperaux}[2]{%
  \IfFileExists{../../build/paper-aux/#2.aux}{%
    \externaldocument[#1]{../../build/paper-aux/#2}%
  }{%
    \IfFileExists{build/paper-aux/#2.aux}{%
      \externaldocument[#1]{build/paper-aux/#2}%
    }{%
      \IfFileExists{#2.aux}{%
        \externaldocument[#1]{#2}%
      }{}%
    }%
  }%
}

% General external reference macros
\newcommand{\paperref}[2]{\ref{#1-#2}}
\newcommand{\papereqref}[2]{(\ref{#1-#2})}

% CPSV16 source (1606.00608 / MPDO-22-12-17-2.tex) external document and shortcuts
\importpaperaux{cpsv16-}{1606.00608/MPDO-22-12-17-2-tnlean}
\newcommand{\cpsvref}[1]{\paperref{cpsv16}{#1}}
\newcommand{\cpsveqref}[1]{\papereqref{cpsv16}{#1}}

% Verdict marker: \gapnote{<kind>}{<status>}. Kind is the policy
% classification (clarification, local-correction, scope-restriction,
% unfaithful, false-source, open-gap); status is open, wip (elimination
% underway), resolved, or historical. Machine-read by the site index;
% prints a verdict line.
\newcommand{\gapnote}[2]{\par\noindent\textbf{Verdict:} #1 (#2).\par}


\title{The Root-Channel Kraus-Rank Step in Theorem 4.1}
\author{}
\date{2026-04-30}

\begin{document}
\maketitle
\gapnote{unfaithful}{open}

\section{The assertion}

This note concerns the reverse implication in Theorem~4.1 of de las Cuevas,
Cirac, Schuch, and P\'erez-Garc\'ia
\cite[Theorem~4.1]{DeLasCuevas2017Irreducible}.\footnote{For traceability, this
note corresponds to
\href{https://github.com/LionSR/TNLean/issues/1046}{issue \#1046}. Related
formal work is recorded in
\href{https://github.com/LionSR/TNLean/issues/670}{issue \#670},
\href{https://github.com/LionSR/TNLean/issues/821}{issue \#821},
\href{https://github.com/LionSR/TNLean/issues/622}{issue \#622}, and
\href{https://github.com/LionSR/TNLean/issues/664}{issue \#664}.} The relevant
source passages are
\path{Papers/1708.00029/main.tex:717--731} and
\path{Papers/1708.00029/main.tex:812--818}.

Let $B=(B^1,\ldots,B^d)$ be a matrix-product-state tensor and let $p\geq 1$.
Its transfer channel is
\begin{align}
    \mathcal E_B(X)
    &=\sum_{i=1}^d B^iX(B^i)^\dagger.
    \label{eq:dccsp17_root_kraus_rank_thm41_transfer_channel}
\end{align}
The source calls a trace-preserving completely positive map $\mathcal E$
$p$-divisible when there is a trace-preserving completely positive map
$\mathcal E'$ such that
\begin{align}
    \mathcal E
    &=(\mathcal E')^p.
    \label{eq:dccsp17_root_kraus_rank_thm41_channel_divisibility}
\end{align}
This condition concerns channels and imposes no bound on the number of Kraus
operators required to represent the root $\mathcal E'$.

The source calls $B$ $p$-refinable when there are a tensor
$A=(A^1,\ldots,A^d)$ on the same physical alphabet and an isometry
\begin{align}
    W
    &:\mathbb C^d\longrightarrow(\mathbb C^d)^{\otimes p}
    \label{eq:dccsp17_root_kraus_rank_thm41_refinement_isometry}
\end{align}
such that
\begin{align}
    \ket{V_{pN}(A)}
    &=W^{\otimes N}\ket{V_N(B)}
    \qquad\text{for every }N.
    \label{eq:dccsp17_root_kraus_rank_thm41_refinement_identity}
\end{align}
Thus the tensor root has exactly the original $d$ physical labels. For $B$ in
irreducible form~II, Theorem~4.1 of
Ref.~\cite{DeLasCuevas2017Irreducible} states the unconditional equivalence
\begin{align}
    B\text{ is }p\text{-refinable}
    \quad&\Longleftrightarrow\quad
    \mathcal E_B\text{ is }p\text{-divisible}.
    \label{eq:dccsp17_root_kraus_rank_thm41_source_equivalence}
\end{align}

\section{The missing Kraus-rank step}

Assume that $\mathcal E_B$ is $p$-divisible. By
(\ref{eq:dccsp17_root_kraus_rank_thm41_channel_divisibility}), there is a
trace-preserving completely positive map $\mathcal E'$ satisfying
\begin{align}
    \mathcal E_B
    &=(\mathcal E')^p.
    \label{eq:dccsp17_root_kraus_rank_thm41_root_channel}
\end{align}
The reverse proof then asserts the existence of a tensor
$A=(A^1,\ldots,A^d)$ for which
\begin{align}
    \mathcal E_B
    &=\mathcal E_A^p.
    \label{eq:dccsp17_root_kraus_rank_thm41_tensor_root}
\end{align}
It proceeds by applying the isometric freedom of Kraus representations to two
representations of the same channel.

The passage from
(\ref{eq:dccsp17_root_kraus_rank_thm41_root_channel}) to
(\ref{eq:dccsp17_root_kraus_rank_thm41_tensor_root}) is not automatic. To
realize $\mathcal E'$ as the transfer channel of a $d$-letter tensor, the root
must have a Kraus representation with at most $d$ terms. Equivalently, its
minimal Kraus rank, which equals its Choi rank, must satisfy
\begin{align}
    \operatorname{rank}_{\mathrm{Kraus}}(\mathcal E')
    &\leq d.
    \label{eq:dccsp17_root_kraus_rank_thm41_required_rank_bound}
\end{align}
If the minimal Kraus rank is $r>d$, then $\mathcal E'$ can be represented by
an $r$-letter tensor but not by a $d$-letter tensor. The resulting isometry
would have codomain $(\mathbb C^r)^{\otimes p}$ rather than the required
codomain $(\mathbb C^d)^{\otimes p}$.

This obstruction is mathematical rather than notational. A completely
positive map can require more than a prescribed number of Kraus operators.
The channel identity
(\ref{eq:dccsp17_root_kraus_rank_thm41_root_channel}) does not exclude the
possibility that every Kraus representation of $\mathcal E'$ has more than
$d$ terms. The reverse implication therefore needs either a proof that the
irreducible-form hypotheses force a root satisfying
(\ref{eq:dccsp17_root_kraus_rank_thm41_required_rank_bound}) or an alternative
argument that avoids a $d$-letter tensor root.

\section{The formal statement}

The formal definition follows the channel-level notion of $p$-divisibility
from Ref.~\cite{DeLasCuevas2017Irreducible}: it asserts the existence of a
trace-preserving completely positive root channel and imposes no Kraus-rank
bound. The conditional reverse theorem separates the missing step from the
rest of the argument.\footnote{The definition is
\leanid{MPSTensor.IsPDivisibleChannel} in
\path{TNLean/MPS/Periodic/Symmetry/Theorem41Defs.lean:24--32}. The conditional
reverse theorem is
\leanid{MPSTensor.thm_4_1_p_refinement_reverse} in
\path{TNLean/MPS/Periodic/Symmetry/Theorem41Reverse.lean:120--142}.}

The additional formal hypothesis is an inverse canonicalization statement. In
ordinary mathematical terms, it says that if $B$ is in irreducible form~II and
$\mathcal E_B$ is $p$-divisible, then there is a $d$-letter tensor $A$ such
that
\begin{align}
    \mathcal E_B
    &=\mathcal E_{A^{[p]}}.
    \label{eq:dccsp17_root_kraus_rank_thm41_inverse_canonicalization}
\end{align}
Given this tensor-level root, the remaining argument follows the source proof.
The blocked tensor $A^{[p]}$ and the original tensor $B$ provide two finite
Kraus representations of the same completely positive map. Isometric freedom
of Kraus representations then gives the isometry
(\ref{eq:dccsp17_root_kraus_rank_thm41_refinement_isometry}) needed for
refinement.\footnote{The inverse canonicalization hypothesis is
\leanid{MPSTensor.PRefinementInverseCanonicalization} in
\path{TNLean/MPS/Periodic/Symmetry/Theorem41Reverse.lean:26--42}. The blueprint
records the bundled conditional equivalence and the conditional reverse
implication in
\path{blueprint/src/chapter/ch22_periodic_ft.tex:2012--2042} and
\path{blueprint/src/chapter/ch22_periodic_ft.tex:2096--2122}.}

The formal development also isolates a sufficient channel-theoretic
hypothesis. If every trace-preserving completely positive root $\mathcal E'$
of $\mathcal E_B$ satisfies
(\ref{eq:dccsp17_root_kraus_rank_thm41_required_rank_bound}), then one may
choose a Kraus representation with at most $d$ terms and pad it with zero
Kraus operators to obtain a $d$-letter tensor $A$. This is a substantive
hypothesis about the root channel. The theorem identifying the least number of
Kraus operators with the Choi rank explains how to express the condition, but
does not prove
(\ref{eq:dccsp17_root_kraus_rank_thm41_required_rank_bound}) for a root of
$\mathcal E_B$.\footnote{The bounded-root formulation is
\leanid{MPSTensor.PRefinementInverseRootKrausRankBound}. Its implication to
inverse canonicalization is
\leanid{MPSTensor.pRefinementInverseCanonicalization_of_rootKrausRankBound},
recorded in
\path{TNLean/MPS/Periodic/Symmetry/Theorem41Reverse.lean:98--118}.}

\section{Conclusion}

The formal reverse theorem has an additional hypothesis relative to
Theorem~4.1 of Ref.~\cite{DeLasCuevas2017Irreducible}. The source claims that
channel-level $p$-divisibility of $\mathcal E_B$ suffices to construct a
$p$-refinement of $B$, but its proof uses a tensor root whose existence
requires the Kraus-rank bound
(\ref{eq:dccsp17_root_kraus_rank_thm41_required_rank_bound}). This bound is
absent from the source definition of $p$-divisibility and is not proved in the
cited converse paragraph.

It remains open whether, under the irreducible-form hypotheses, a channel root
can always be chosen with Kraus rank at most $d$. The formal development
therefore proves the reverse implication under the explicit inverse
canonicalization hypothesis
(\ref{eq:dccsp17_root_kraus_rank_thm41_inverse_canonicalization}), or under the
stronger root-Kraus-rank assumption that implies it. Until a paper-level
argument establishes the required rank bound, the reverse implication must
remain conditional on this additional mathematical input.

\bibliographystyle{alpha}
\bibliography{references}

\end{document}
